\section{Relative dynamics and proximity operations}
\label{sec:proximity}

\subsection{The injection error is the whole point}

The rendezvous of Section~\ref{sec:transfer} closes to a fraction of a
millimetre, which no ascent stage has ever achieved. The dispersion listed in
Section~1.1 -- single-string propulsion, thrust-to-centre-of-mass misalignment,
active-only attitude control, migrating inertia -- arrives at the top of the
climb as a position and velocity offset. This part models that offset explicitly
and prices its removal.

The error is drawn once from a seeded generator: a random direction with
magnitude uniform in $[500, 1000]\unit{m}$ for position and $[0.1, 1]\unit{m/s}$
for velocity. The realisation used throughout is
$\lVert\delta\mathbf r_0\rVert = \mDrZero\unit{m}$ and
$\lVert\delta\mathbf v_0\rVert = \mDvZero\unit{m/s}$. The seed, not the numbers,
is the reproducibility mechanism.

\subsection{Frame definition, and the sign trap inside it}
\label{sec:proximity:frame}

Relative motion is expressed in a local frame attached to the target. This
report uses

\begin{equation}
\hat{\mathbf x} = -\hat{\mathbf v}, \qquad
\hat{\mathbf y} = -\hat{\mathbf h}, \qquad
\hat{\mathbf z} = +\hat{\mathbf r},
\label{eq:lvlh}
\end{equation}

that is, $x$ along the \vbar{} measured positive \emph{astern} of the target,
$y$ cross-track along the negative orbit normal, and $z$ along the \rbar{},
positive outward. The triad is right-handed, since
$\hat{\mathbf x}\times\hat{\mathbf y} = (-\hat{\mathbf v})\times(-\hat{\mathbf h})
= \hat{\mathbf v}\times\hat{\mathbf h} = \hat{\mathbf r}$.

\begin{quote}\small\itshape
\textbf{This is not Fehse's frame.} Fehse's $+$\vbar{} points ahead of the
target; here it points astern. A ``$50\unit{m}$ \vbar{} hold'' in this report is
therefore $50\unit{m}$ \emph{behind} the target, and the sign of every Coriolis
term below is flipped relative to Fehse or to any radial/along-track
formulation. Every number in Part~3 is internally consistent with this choice.
They are not interchangeable with a Fehse implementation without negating $x$
and $\dot x$. Reference~\cite{fehse} is cited for the linearisation and for the
operational shape of the approach, not as the source of this axis convention.
\end{quote}

Two remarks justify what looks at first like a perverse choice. Operationally,
\vbar{} range is quoted astern: a ``$50\unit{m}$ \vbar{} hold'' is $50\unit{m}$
behind the target, and the approach then walks $x$ from $+50\unit{m}$ down to
zero at the port. Mathematically, the price is that the target's angular
velocity is $-n$ about the frame's own $y$ axis, which flips the sign of every
Coriolis term relative to the textbook radial/along-track ordering. Linearising
about a circular chief \cite{cw1960} gives

\begin{align}
\ddot x - 2n\dot z &= 0, \label{eq:hcw1}\\
\ddot y + n^2 y &= 0, \label{eq:hcw2}\\
\ddot z + 2n\dot x - 3n^2 z &= 0. \label{eq:hcw3}
\end{align}

This is exactly where a sign error hides most comfortably: copy a state
transition matrix written for $x$ radial, label the axes \vbar{} and \rbar{}, and
the free-drift plot still looks plausible while every targeting solution is
quietly wrong. The convention was therefore settled empirically rather than
asserted, and it is now pinned by \texttt{tests/test\_hcw\_convention.m}, which
checks $\Phi(0) = I$, $\Phi(t)\Phi(-t) = I$, $\dot\Phi(0) = A$, the closed form
against a column-by-column numerical integration of the same linear system to
$3\times10^{-9}$, the sign and magnitude of the radial-release drift, and the
reference cross-check below.

Fed an independent implementation's injection error, the targeting of
Section~\ref{sec:proximity:hold} returns $1.419\unit{m/s}$ against that
implementation's $1.4\unit{m/s}$. Reading the same physical state in the
ahead-positive frame instead returns $1.236\unit{m/s}$ and does not reproduce
it. That is the evidence for the frame, and it is also why the ahead-positive
convention was not adopted: switching would have broken the only independent
check Part~3 has.

A physical check completes the argument. Release the chaser $100\unit{m}$
radially outward at zero relative velocity: it is on a higher, slower orbit and
must fall behind. The state transition matrix returns $x = +11.3\unit{km}$ after
three orbits, and since $+x$ is astern, it has indeed fallen behind. Reading
$+x$ as ``ahead'' is the trap; the sign is correct once the axis direction is
stated.

\subsection{State transition matrix}

With $\tau = nt$ and the state $\delta\mathbf x = [x\ y\ z\ \dot x\ \dot y\ \dot z]^\top$,
equations~\eqref{eq:hcw1}--\eqref{eq:hcw3} integrate in closed form to
$\delta\mathbf x(t) = \Phi(t)\,\delta\mathbf x(0)$ with

\begin{equation}
\Phi(t) =
\begin{bmatrix}
1 & 0 & 6(\tau-\sin\tau) & \frac{4\sin\tau-3\tau}{n} & 0 & \frac{2(1-\cos\tau)}{n}\\[2pt]
0 & \cos\tau & 0 & 0 & \frac{\sin\tau}{n} & 0\\[2pt]
0 & 0 & 4-3\cos\tau & \frac{2(\cos\tau-1)}{n} & 0 & \frac{\sin\tau}{n}\\[2pt]
0 & 0 & 6n(1-\cos\tau) & 4\cos\tau-3 & 0 & 2\sin\tau\\[2pt]
0 & -n\sin\tau & 0 & 0 & \cos\tau & 0\\[2pt]
0 & 0 & 3n\sin\tau & -2\sin\tau & 0 & \cos\tau
\end{bmatrix}.
\label{eq:stm}
\end{equation}

Three properties are checked automatically: $\Phi(0) = I$ to machine precision;
$\dot\Phi(0)$ equals the system matrix of the first-order form, which is what
actually pins the Coriolis signs; and $\Phi(t)\Phi(-t) = I$.

\subsection{Free drift}

Propagating the injection error with no correction is instructive.
Figure~\ref{fig:drift} shows the result over three target orbits: the
out-of-plane component is a bounded harmonic at the orbital frequency, while the
in-plane motion develops a secular along-track drift driven by the period
mismatch. The separation reaches $\mDriftKm\unit{km}$.

That number is the argument for active proximity operations in one line. A
$0.88\unit{m/s}$ velocity error, roughly $0.05\%$ of orbital speed, becomes a
$21\unit{km}$ miss within seven hours.

\begin{figure}[htbp]
\centering
\includegraphics[width=0.78\textwidth]{fig06_free_drift_3d.png}
\caption{Uncorrected injection error over three target orbits in the local
frame. Starting from $\mDrZero\unit{m}$ and $\mDvZero\unit{m/s}$, the chaser
drifts $\mDriftKm\unit{km}$ along-track. Cross-track motion stays bounded.}
\label{fig:drift}
\end{figure}

\subsection{Two-impulse hold acquisition}
\label{sec:proximity:hold}

Partitioning~\eqref{eq:stm} into $3\times3$ blocks, the transfer from the
injection state to a prescribed hold point in a prescribed time is a single
linear solve:

\begin{equation}
\mathbf v_0^+ = \Phi_{rv}^{-1}\!\left(\mathbf r_{\text{hold}} - \Phi_{rr}\,\delta\mathbf r_0\right),
\qquad
\dv_1 = \mathbf v_0^+ - \delta\mathbf v_0,
\label{eq:target1}
\end{equation}
\begin{equation}
\dv_2 = \mathbf v_{\text{hold}} - \left(\Phi_{vr}\,\delta\mathbf r_0 + \Phi_{vv}\,\mathbf v_0^+\right).
\label{eq:target2}
\end{equation}

The hold point is $\mathbf r_{\text{hold}} = [50, 0, 0]^\top\unit{m}$ and the
transfer time $0.3\,T_{MS} = \mDtTr\unit{s}$. The cost is
$\dv_{\text{hold}} = \mDVhold\unit{m/s}$, and propagating the resulting
two-burn trajectory back through~\eqref{eq:stm} arrives within
$10^{-13}\unit{m}$ of the target point.

The block $\Phi_{rv}$ becomes singular whenever the transfer time approaches a
multiple of the orbital period, because after a full revolution the initial
velocity no longer controls the final position. Its condition number is
therefore checked before every solve and the transfer time nudged by $1\%$ until
the inversion is safe; at $0.3\,T_{MS}$ the reciprocal condition number is
$0.21$, comfortably clear.

\subsection{Forced straight-line \vbar{} approach}

From the hold point the vehicle is walked to the port in $N = 5$ equal legs of
$200\unit{s}$, each targeting the waypoint
$\mathbf r_k = \mathbf r_{\text{hold}}(1 - k/N)$, followed by a braking impulse
that nulls the residual velocity at contact. The total is
$\dv_{\text{dock}} = \mDVdock\unit{m/s}$, and the trajectory terminates
$3\times10^{-16}\unit{m}$ from the origin.

This profile is not chosen for elegance. A sequence of short targeted hops keeps
plume impingement away from the target, bounds the closing rate between
checkpoints, and gives the crew or the ground a natural abort opportunity at
every waypoint -- the pattern used on real \vbar{} approaches \cite{fehse}.

\begin{figure}[htbp]
\centering
\includegraphics[width=0.93\textwidth]{fig08_forced_vbar_docking.png}
\caption{Forced straight-line approach. Top: \vbar{} and \rbar{} components
against time, with the six impulses marked. Bottom: the in-plane path from the
$50\unit{m}$ hold to the docking port, total cost $\mDVdock\unit{m/s}$.}
\label{fig:docking}
\end{figure}

The full proximity phase therefore costs $\mDVprox\unit{m/s}$, about $1\%$ of
the $\mDVtot\unit{m/s}$ transfer budget. That is the price of the ascent
dispersion, and it is only obtainable by modelling both phases in the same
units.

\subsection{Verification against the nonlinear equations}

The linear model is finally checked against the exact nonlinear relative
equations of motion for a circular chief, integrated with the same tolerances.
Writing $\rho^2 = x^2+y^2+z^2$ and $r_c = \sqrt{x^2+y^2+(R_2+z)^2}$, and using
the regularised form

\begin{equation}
q = \frac{\rho^2 + 2R_2 z}{R_2^2}, \qquad
F = \frac{q\left(2 + q + \sqrt{1+q}\right)}{(1+q)^{3/2}\left(1+\sqrt{1+q}\right)},
\label{eq:nonlin}
\end{equation}

avoids differencing two nearly equal inverse cubes, which matters when
$\rho/R_2 \sim 10^{-5}$.

The two models part company by $\mNlSep\unit{m}$ after three orbits at kilometre
range -- about half a percent of the separation. At the $50\unit{m}$ hold point,
where $\rho/R_2 = 2.3\times10^{-5}$, they agree to well under a millimetre. The
linearisation is used exactly where it is valid, and the divergence at long
range is quantified rather than assumed away.
